\section{A Checked and Modular Proof Workflow}
\label{sec:proof-workflow}
\label{sec:case-study}

When automated reasoning leaves obligations open, completing the development requires manual interaction with the prover.
The user may then introduce partial expressions absent from the specification, establish intermediate claims, or apply lemmas proved elsewhere.
The proof environment must check the additional preconditions and preserve the assumptions on which reused facts depend.
We compare how Atelier~B\footnote{Community Edition 24.04.2.} and \barrel support this work, focusing on well-definedness during interaction, correctness of inference rules, reusable mathematical lemmas and explicit proof dependencies.

The examples draw on the six-component B development of the IEEE~1394 tree identify protocol by Abrial, Cansell and M\'ery~\cite{AbrialCansellMery2003}, which refines a leader-election algorithm.\footnote{We thank Dominique M\'ery for giving us access to the B refinement chain and for the insightful discussions that helped us understand it.}
We ran Atelier~B's Force~3 followed by its predicate prover (\texttt{pp}) on the generated obligations, then completed the remaining proofs interactively.
The artifact contains both completed developments and the counterexamples and dependency checks below.

\subsection{Well-definedness throughout interaction}
\label{sec:lean-ip-comparison}

The WD obligations generated from a B specification concern the expressions occurring in that specification.
Interactive proof can introduce further partial expressions, for instance when choosing an existential witness.
Their definedness must be justified at that point.

Atelier~B's reference manual warns that the tactic \texttt{se}, which supplies an existential witness, does not protect against ill-typing or ill-definedness~\cite[\S4.44]{AtelierBInteractiveProver}.
Consider a machine containing a contradictory assertion:
\begin{lstlisting}[language=b, basicstyle=\small\ttfamily, numbers=none]
MACHINE DerivFalse
ASSERTIONS #xx.(xx : INTEGER & xx /: INTEGER)
END
\end{lstlisting}
After an initial deduction step (\texttt{dd}), the command \inlineab!se(min({}))! selects the minimum of the empty set as witness.
The displayed goal becomes
\[
  \bmin{\varnothing} \in \mathbb{Z}
  \;\land\;
  \bmin{\varnothing} \notin \mathbb{Z}.
\]
A call to \texttt{mp}, the mini-prover, closes it.
The prover therefore accepts a contradictory assertion after an undefined witness has been supplied interactively.

In \barrel, forming \(\min S\) requires evidence that $S$ has a least element (\cref{lst:min-fn-app}).
The witness \(\min\varnothing\) thus requires a proof of \(\mathtt{min.WD}\ \varnothing\), that is,
\[
  \exists x \in \varnothing,\ \forall y \in \varnothing,\ x \le y,
\]
which is unprovable.
Admitting this WD condition triggers a Lean warning, and the \texttt{assert\_no\_sorry} check in \texttt{examples/DerivFalse.lean} fails on the resulting theorem.
For operators encoded with WD parameters, definedness is therefore checked throughout proof construction, including steps the source PO generator could not anticipate.

\subsection{Trustworthy and modular library of lemmas}
\label{sec:checked-rules}

Definedness checks cannot compensate for an incorrect inference rule.
Among Atelier~B's thousands of shipped rules, \texttt{GenPredicateX.8} gives the following invalid inference:
\[
  \frac{a \notin \mathcal{P}(X) \cup \{y\}}
       {\exists z\in a,\ z\notin X \land z\ne y}.
\]
The premise places $y$ at the type of $a$, whereas the conclusion compares it with an element of $a$.
The schema therefore has no consistent B typing for arbitrary \(y\).
Instantiating $y$ with the polymorphic empty set nevertheless gives individually well-typed formulas and an invalid inference.

Indeed, let $a := \{\varnothing\}$ and $X := \mathcal{P}_1(\mathbb{Z})$, the set of nonempty subsets of $\mathbb{Z}$.
Since $\varnothing \in a$ but $\varnothing \notin X$, we have $a \notin \mathcal{P}(X)$. Since $a \ne \varnothing$,
the premise $a\notin\mathcal{P}(X)\cup\{\varnothing\}$ holds.
The conclusion requires an element of $\{\varnothing\}$ distinct from $\varnothing$, which is impossible.

The following script allows one to close any proof using this rule, independently of any WD obligation:
\begin{lstlisting}[language=b, basicstyle=\small\ttfamily, numbers=none, mathescape=false, escapeinside={}]
ah(not({{}}: POW(POW1(INTEGER)) \/ {{}})) & pr & dd &
  ar(GenPredicateX.8,Fwd) & pr
\end{lstlisting}

In Lean, we refute this rule instance and prove that replacing \(\mathcal{P}(X) \cup\{y\}\) by $\mathcal{P}(X\cup\{y\})$ in the premise yields a valid rule.
This is available in the artifact in the file \texttt{examples/Falsehood.lean}.

Users can extend \barrel with their own lemmas, drawing on Mathlib to construct proofs that Lean's kernel checks.
The supplied library and user extensions therefore require no separate trust in the correctness of their statements: each lemma is justified by its proof.

\begin{example}[Definedness after a function update]
Updating a total function at a point of its domain preserves totality when the new value belongs to its codomain:
\begin{equation}
  {f \in A \to B} \;\land\; {x \in A} \;\land\; {y \in B}
  \quad\implies\quad
  f \bovr \{x \mapsto y\} \,\in\, A \to B.
  \label{eq:workflow-override}
\end{equation}
Atelier~B's related rule \texttt{InFunctionXY.22} establishes preservation of partial-function typing under overriding.
It is a special case of \barrel's proved lemma \texttt{pfun\_of\_overload}; totality also requires a domain argument.
The lemma \texttt{tfun\_\allowbreak override\_\allowbreak singleton} proves \cref{eq:workflow-override} for arbitrary carrier types using Mathlib's sets and relations.
Combining it with the WD lemma for total-function application gives a theorem that can be used in any obligation involving such an update:
\begin{lstlisting}[language=lean, basicstyle=\small\ttfamily, numbers=none]
theorem update_wd {α β : Type}
    {A : Set α} {B : Set β} {f : SetRel α β}
    {x z : α} {y : β}
    (hf : f ∈$$ A ⟶$$ B) (hx : x ∈$$ A)
    (hy : y ∈$$ B) (hz : z ∈$$ A) :
    app.WD (f <+ {(x, y)}) z :=
  app.WD_of_mem_tfun (tfun_override_singleton hf hx hy) hz
\end{lstlisting}
\end{example}

In our Atelier~B scripts, intermediate claims are proved in the context of the current PO, and recurring arguments are repeated at their uses.
Atelier~B also supports assertions and user rules, but applying a rule does not require a kernel-checked proof of its correctness.
The Lean approach combines reuse with a check of the general lemma, using the same proof language and library as the obligations themselves.

\subsection{Refinement with explicit proof dependencies}

The WD parameters of encoded operators must also be supplied in expressions inherited from an abstract component.
When that component has already been proved in Lean, \barrel can supply the evidence by applying its WD theorems.
Lean checks these applications in the refinement's context, so reuse avoids a new proof obligation without omitting the definedness check.

After \texttt{qed}, a component's WD theorems are available to subsequent imports, including across Lean modules.
Where possible, \barrel removes unnecessary premises from their reusable statements, allowing later components to apply them with fewer hypotheses.
If no available theorem applies, \barrel generates a new WD obligation.

\begin{example}[Reusing definedness across refinements]
The refinement \texttt{mm2} inherits an assertion from \texttt{mm1} containing $\mathtt{ff}(x)$, where \texttt{ff} assigns a tree to each node in \texttt{ND}.
Definedness of this application needs only
\[
  \mathtt{ff} \in \mathtt{ND} \to (\mathtt{ND} \pfun \mathtt{ND})
  \quad\land\quad x \in \mathtt{ND}
  \quad\Longrightarrow\quad \mathtt{app.WD}\ \mathtt{ff}\ x.
\]
Atelier~B generates WD obligations for the assertion in \texttt{mm1} and does not generate them again in \texttt{mm2}.
In \barrel, importing \texttt{MM1.lean} into \texttt{MM2.lean} makes the proved WD fact available when encoding the refinement.
All 53 occurrences of the inherited application then share evidence obtained from \texttt{mm1}'s theorem, with no new WD obligation in \texttt{mm2}.
The proof terms contain the application of the earlier theorem, which Lean checks using the hypotheses available in \texttt{mm2}.
\end{example}

Across the six components, Atelier~B generates 60 WD obligations from the B specification.
\barrel instead extracts WD conditions while translating individual expressions and shares their proofs between occurrences.
It generates 31 WD goals, all proved automatically, and avoids 7 further goals by reusing three previously proved theorems across components.
These goal counts exclude repeated uses of the same evidence, such as the 53 applications in \texttt{mm2}; every encoded partial application still carries a WD proof.

\subsection{Automation, interaction and trust}

Atelier~B's automation, developed specifically for B over more than twenty years, leaves fewer obligations for interactive proof in this development.
The two counterexamples nevertheless show that a successful prover status can conceal an invalid inference.
In \barrel, Lean's kernel checks the resulting proof terms and the proofs of every lemma they use.

Interactive proof can also help extend trustworthy automation in \barrel.
Registering the totality lemma in \cref{eq:workflow-override} with the tactic for total-function membership lets the WD tactic combine it with domain membership, as in the explicit proof above.
With these generic lemmas and extended premise search, \barrel automatically discharges all locally generated WD goals in the leader election development.
Adding a tactic or lemma does not add a trusted inference rule: its use must produce evidence accepted by Lean's kernel.

The proof environment also affects the effort needed when automation fails.
Atelier~B displays goals, hypotheses and proof histories, but an output message like ``The Predicate Prover failed to prove the current goal'' does not distinguish missing premises from unsuccessful search.
Lean lets the user isolate a claim with \texttt{have}; applying a theorem exposes its remaining premises as subgoals, and a failed \texttt{exact} reports the supplied and expected types.
These diagnostics provide a concrete point at which to investigate the failed step, although matching and normalization can still require manual work.

Using \barrel requires B users to learn Lean and the relevant Mathlib abstractions, which may be a barrier to adoption. In return, Lean checks the proofs produced by automation and the lemmas they use, including explicit WD evidence across components. The counterexamples above show why this protection matters even when a specialized prover offers stronger automation.